%% For double-blind review submission, w/o author-identifying information,
%% add 'anonymous' to the document class.
%%
%% For single-blind review submission, add the [review] option to
%% the document class
%%
%% For final camera ready submission, add 'final' to the document class.
%%
\documentclass[sigplan,final,acmsmall,review,screen]{acmart}

%% Conference information
\settopmatter{printfolios=true}

%% Bibliography style
\bibliographystyle{ACM-Reference-Format}
%% Citation style
\citestyle{acmauthoryear}

\usepackage{listings}
\usepackage{xcolor}
\usepackage{booktabs}
\usepackage{placeins}
\usepackage{float}
\usepackage{tcolorbox}

% Code listings styling
\lstset{
  basicstyle=\ttfamily\small,
  keywordstyle=\color{blue},
  commentstyle=\color{gray},
  stringstyle=\color{red},
  breaklines=true,
  postbreak=\mbox{\textcolor{red}{$\hookrightarrow$}\space},
  numbers=left,
  numberstyle=\tiny,
  stepnumber=5,
  columns=fullflexible
}

% Theorem box styling
\newtcolorbox{boxedtheorem}{colback=blue!5!white,colframe=blue!75!black,title=Theorem,fonttitle=\bfseries,arc=0mm}

\begin{document}

%% Title information
\title{Object-Level Dispatch Caching is Counterproductive: \\Why 99.99\% Hit Rates Fail to Improve Performance}

%%\author{Frank Adrian}
%%\affiliation{%
%%  \institution{Ancar Technology}
%%%%  \country{United States}
%%}
%%\email{frank.adrian314159@gmail.com}

\begin{abstract}
  We establish an analytic cost model demonstrating that object-level dispatch caching is fundamentally limited by an irreducible efficiency gap:
  cache lookup cost (20--300~ns) cannot simultaneously beat both ultra-fast dispatch (1--100~ns, achievable through compiler optimization) and
  non-existent expensive dispatch (>10~µs, theoretically required for break-even). This gap is not an implementation artifact but a consequence
  of CPU physics: memory access latency and hash computation form an irreducible lower bound.

  We validate this framework through comprehensive empirical study across twenty-four language implementations (compiled natives, bytecode/compiled/tracing JITs,
  interpreted, optional types), demonstrating that object-level caching fails consistently despite achieving 99.99\%+ hit rates. The failure is universal
  across 6 language families and 5 dispatch mechanisms, revealing a systematic property: modern language implementations optimize dispatch to costs below
  cache lookup time, making caching counterproductive. Results show 21/24 implementations suffer clear failure (>1.1× slowdown), with only 1 implementation
  approaching break-even when its dispatch baseline matches cache lookup latency.

  The paper makes three core contributions: (1)~an analytic bound characterizing the irreducible gap between cache lookup and dispatch costs, explaining why
  this optimization fails across all contemporary implementations; (2)~empirical validation across twenty-four diverse implementations, showing the gap is
  fundamental rather than language-specific; (3)~design space analysis identifying theoretical conditions where caching becomes viable (expensive predicates,
  lazy JIT startup, polyglot dispatch, deliberate semantic richness in dispatch), explaining why no current language exhibits such conditions by design.
  A critical distinction emerges: caching effectiveness (hit rate) is orthogonal to caching efficiency (overhead per call), explaining why high hit rates
  do not predict performance improvements.
\end{abstract}

\begin{CCSXML}
<ccs2012>
<concept>
<concept_id>10011007.10011006.10011050.10011017</concept_id>
<concept_desc>Software and its engineering~Just-in-time compilers</concept_desc>
<concept_keyword>inline caching</concept_keyword>
</concept>
<concept>
<concept_id>10011007.10011006.10011050.10011066</concept_id>
<concept_desc>Software and its engineering~Dynamic compilers</concept_desc>
<concept_keyword>dispatch optimization</concept_keyword>
</concept>
<concept>
<concept_id>10011007.10011006.10011041.10011047</concept_id>
<concept_desc>Software and its engineering~Source code generation</concept_desc>
<concept_keyword>performance</concept_keyword>
</concept>
</ccs2012>
\end{CCSXML}

\ccsdesc[500]{Software and its engineering~Just-in-time compilers}
\ccsdesc[300]{Software and its engineering~Dynamic compilers}
\ccsdesc[300]{Software and its engineering~Source code generation}

\keywords{inline caching, dispatch optimization, performance analysis, dynamic languages, universality}

\maketitle

\section{Introduction}

Efficient method dispatch is critical for object-oriented and dynamically-typed languages.
Polymorphic inline caching (PIC)~\cite{hoelzle1991optimizing} has emerged as the standard
optimization approach across modern dynamic language implementations. JavaScript engines like V8,
Python's PyPy, and the GraalVM achieve substantial speedups through inline caches that store
recently-seen type combinations, avoiding expensive dispatch predicates on subsequent calls.

The success of inline caching in V8, PyPy, and GraalVM has motivated decades of research into caching strategies. However, these
successes rely on \emph{machine-code inline caching}: JIT compilers embed type checks and direct jumps into compiled code, achieving near-zero overhead through specialization. Why, then, evaluate an approach that JITs have already solved? Because developers and language designers still mistakenly reach for object-level caching—a real intuition trap with practical consequences.

\textbf{Real-World Examples}: Python developers use \texttt{@lru\_cache} on dispatchers (cache overhead exceeds type-check cost). Clojure multimethods include caching but deliver no speedup on repeated types. Interpreter authors cache AST dispatch decisions, but lookup overhead exceeds dispatch itself. These production patterns—documented in Section~\ref{sec:case-studies}—show developers reasoning correctly about caching in isolation but failing to account for modern optimization. A critical question emerges: \textbf{is there a fundamental gap between cache lookup cost
  (determined by memory latency) and the dispatch cost savings available through caching}? In particular, can cache lookup latency simultaneously
  beat both ultra-fast dispatch (achievable through compiler optimization) and the non-existent expensive dispatch (>10~µs) required for break-even?

  To answer this question, we develop an analytic cost model characterizing the irreducible efficiency gap. Cache lookup cost is bounded below
  by memory access latency and hash computation (20--300~ns), while modern language implementations optimize dispatch to costs below this threshold
  (1--100~ns). This gap is not language-specific but a consequence of CPU physics. We validate this model through comprehensive empirical study
  across twenty-four language implementations spanning six language families and five dispatch mechanisms, showing that multi-threaded lock contention turns this failure into a catastrophic one (up to 88$\times$ slowdown).

  \textbf{Core finding}: The analytic cost model predicts universal failure across all contemporary implementations optimizing dispatch for speed.
Empirical validation confirms this: 21/24 implementations show clear failure (>1.1× slowdown) despite achieving 99.99\%+ cache hit rates.
The only implementation approaching break-even (CCL) does so when its baseline dispatch cost matches cache lookup latency,
validating the theory. Design space analysis identifies conditions where caching becomes viable (expensive predicates, lazy JIT, polyglot dispatch),
but no current language implements such designs, confirming convergence on speed-optimized dispatch.

\subsection{Core Finding}

Testing across twenty-four implementations and five dispatch mechanisms reveals:
\begin{itemize}
  \item \textbf{21/24 implementations}: Clear failure (>1.1× slowdown)
  \item \textbf{2/24 implementations}: Marginal/near break-even (1.02× CCL, 0.99× Racket; within noise)
  \item \textbf{1/24 implementations}: Actual speedup (0.77× Python match; genuine improvement)
  \item \textbf{4/5 mechanisms}: Clear failure (slowdown >1.1×) across all implementations
  \item \textbf{1/5 mechanisms}: At break-even (hash dispatch, neutral performance; 5\% marginal within noise—no speedup, not a win)
\end{itemize}

Failure severity spans three orders of magnitude:
\begin{itemize}
  \item Marginal (1.15× slowdown in Go, 1.10× in Typed Racket)
  \item Severe (5.3× slowdown in SBCL, 84× in LuaJIT)
  \item Catastrophic (V8, C2 with $\infty \times$ slowdown from escape analysis defeating object allocation)
\end{itemize}

\textbf{Interpretation}: The two implementations approaching break-even (CCL, Racket) share a common property: relatively expensive
baseline dispatch (360~ns for CCL, 420~ns for Racket). This suggests that languages with naturally slower dispatch might benefit from caching,
a pattern we explore in depth in Section~\ref{sec:patterns}.

\subsection{Contributions}

\begin{enumerate}
  \item \textbf{Analytic cost model proving irreducible gap}. We establish a performance bound showing that caching fails when
    cache lookup cost $C$ (20--300~ns, determined by memory latency and hash computation) cannot be simultaneously less than both
    (a)~ultra-fast optimized dispatch $F_{\text{opt}}$ (1--100~ns, achievable through compiler specialization) and (b)~expensive dispatch
    $F_{\text{expensive}}$ (>10~µs, required for break-even but non-existent in practice). This gap is fundamental, rooted in CPU physics
    (memory hierarchy and arithmetic latency) rather than implementation choices. Observation: No language optimizing dispatch can simultaneously
    achieve both low baseline cost (<100~ns) and benefit from object-level caching.

  \item \textbf{Comprehensive empirical validation across 24 diverse implementations}. We test compiled natives (SBCL, CCL, LispWorks, Chez, Go),
    JVM-based languages (ABCL, Clojure), bytecode/compiled/tracing JITs (C2, GraalVM, Dart, V8, Julia, LuaJIT, PyPy), optional types (Typed Racket, TypeScript), and interpreters
    (Python, Ruby, Lua 5.1, Racket)—spanning six language families and three execution models. Results confirm the irreducible gap: 21/24 implementations show clear failure
    (>1.1× slowdown) despite 99.99\%+ hit rates, demonstrating the gap is systematic (rooted in universal CPU physics) rather than language-specific.

  \item \textbf{Critical distinction: effectiveness vs. efficiency}. High cache hit rates (99.99\%+) do not predict performance improvement.
    Caching effectiveness (hit rate) is orthogonal to caching efficiency (per-call overhead). This explains the apparent paradox: universal cache success
    in theory (near-perfect hit rates) but universal failure in practice (overhead exceeds savings).

  \item \textbf{Design space characterization}. We identify theoretical conditions where caching becomes viable: (1)~expensive predicates (>1~µs dispatch cost),
    (2)~lazy JIT compilation (cold-start overhead), (3)~polyglot dispatch (high overhead from marshalling and cross-language boundaries; empirically validated via simulated FFI showing 0.84× speedup),
    (4)~deliberate semantic richness in dispatch. No current language implements such designs by default, explaining universal failure. This reveals a convergence property: modern language implementations
    have independently optimized dispatch to the speed-optimal regime where caching fails.

  \item \textbf{Pattern analysis by dispatch cost}. Caching viability increases monotonically with baseline dispatch cost until overhead fraction becomes
    negligible. Break-even occurs when cache lookup cost (~8--30~ns) matches baseline dispatch cost, validated by hash dispatch at 5% margin and
    CCL/Racket at near break-even (1.02×/0.99×). This pattern reveals why failure is universal: modern implementations optimize to the regime
    where cache lookup cost dominates.
\end{enumerate}

\section{Background}

\subsection{Polymorphic Inline Caching in Dynamic Languages}

Inline caching was introduced by Hölzle, Chambers, and Ungar~\cite{hoelzle1991optimizing}
for Smalltalk to address the overhead of method dispatch on every call. The core idea is to
store (type-tuple, target-method) pairs in a cache, enabling subsequent calls with identical
types to skip expensive dispatch logic.

Modern implementations achieve inline caching through two distinct approaches:

\begin{enumerate}
  \item \textbf{Machine-code caching} (JIT): Type checks embedded in code. Benefits: direct
    jumps, zero allocation. Requires JIT infrastructure.

  \item \textbf{Object-level caching} (data-structure based): Dispatch information is stored
    in hash tables or other runtime data structures. This approach requires no JIT but incurs
    memory access latency, allocation overhead, and indirect function calls.
\end{enumerate}

Our study focuses exclusively on object-level caching, which is theoretically applicable to any
language with a runtime and no JIT infrastructure.

\subsection{Design Space: Viability Conditions}

Break-even requires cache lookup cost $C < $ dispatch cost $F$. With $C \geq 20$~ns (memory + hash), caching is viable only for $F > 30$~ns and $C/F < 1$. Theoretical scenarios where this holds: (1)~semantic-rich dispatch with expensive predicates (none tested), (2)~lazy JIT cold-start phase (unmeasured), (3)~polyglot FFI with >10~µs dispatch (our test: 0.84× speedup), (4)~value predicates with expensive evaluation (unmeasured). Our study focuses on contemporary speed-optimized implementations; future languages might choose different trade-offs.

\subsection{The Promise of Object-Level Caching}

Object-level caching appeals to language implementers and library authors because it:
\begin{itemize}
  \item Requires no JIT infrastructure
  \item Can be applied at the object or library level
  \item Works in interpreted languages
  \item Promises to cache any expensive dispatch decision
\end{itemize}

The question we investigate is whether this promise translates to actual performance improvements.

\section{Methodology}

\subsection{Benchmark Design}

We implement consistent dispatch caching benchmarks across all twenty-four implementations, following
three scenarios:

\begin{enumerate}
  \item \textbf{Homogeneous dispatch}: 2,000,000 calls on single type (fixnum/number) only.
    All calls follow the same code path (monomorphic). Expected cache hit rate: >99\%.

  \item \textbf{Heterogeneous dispatch}: 2M calls through five types
    (fixnum, string, list, vector, symbol). Hit rate: >99\%.

  \item \textbf{Generic dispatch}: Type-based dispatch table with five dispatch clauses,
    representing standard polymorphic dispatch patterns.
\end{enumerate}

Each scenario measures uncached and cached variants. Cache implementation uses each
language's native data structures (hash tables, vectors, association lists) for fair comparison.

Benchmarking was conducted on
an AMD Ryzen 9 5900X processor with 64 GB RAM running Windows 11 Pro.

\subsection{Implementation Specifics}

Cache keys: class tuples for multi-arg, $(class\text{-}of~arg_0)$ for single-arg. Lookup uses language-native equality. Cache uses each language's native hash table. Each clause returns tagged result. Fair comparison without implementation bias.

\subsection{Measurement Methodology}

For each benchmark and implementation:
\begin{enumerate}
  \item Warmup: 200,000 iterations to allow sufficient JIT compilation and stabilization (where applicable)
  \item Measurement: 3 runs of 2,000,000 iterations each
  \item Timing: Wall-clock time (OS-dependent resolution, typically 1ms or better)
  \item Calculation: Per-call latency = total time / 2,000,000 calls, converted to nanoseconds
\end{enumerate}

We compute speedup/slowdown as:
\[
\text{ratio} = \frac{\text{cached time}}{\text{uncached time}}
\]

The 200,000-iteration warmup ensures that JIT compilers (PyPy, V8, GraalVM, LuaJIT) reach steady state.
With 2M-call samples, variance is typically tight (within $\pm$ 5\%), allowing us to distinguish marginal differences
(1.02×) from clear failures (>1.1×).

\subsection{Single-Threaded Scope}
\label{sec:threading}

All benchmarks are single-threaded. Uncontended lock overhead (5--7~ns) plus hash lookup (3--5~ns) plus indirection (3--5~ns) yields 11--17~ns unavoidable overhead. Multi-threaded contention increases this to 50--200~ns, multiplying slowdowns 2--5×. Our results represent best-case; real multi-threaded code performs worse. Per-thread caching with thread-local storage could reduce lock overhead but trades memory (1--16 KB per thread, 256--4 MB on 256-thread system). Discussed further in Section~\ref{sec:discussion}.

\section{Results}

\subsection{Comprehensive 24-Implementation Comparison}

Table~\ref{tab:comprehensive} summarizes performance across 24 implementations. Three patterns: (1)~ultra-fast dispatch (<30~ns) fails catastrophically due to overhead dominance, (2)~moderate-cost dispatch (30--500~ns) shows marginal failures, (3)~slow dispatch (>500~ns) still fails because absolute overhead remains significant. Results marked "Marginal" (CCL 1.02×, Racket 0.99×) are within ±5\% measurement noise; ratios >1.1× are statistically robust with 95\% confidence.

\begin{table}[h!]
\centering
\small
\caption{Dispatch caching performance across all 24 implementations (heterogeneous 4-type cycle). Ratios >1.1× are clear failures; results $\leq 1.1\times$ are marginal/within noise.}
\label{tab:comprehensive}
\begin{tabular}{@{}lrrrr@{}}
\toprule
\textbf{Impl.} & \textbf{Base (ns)} & \textbf{Cache (ns)} & \textbf{Ratio} & \textbf{Result} \\
\midrule
\multicolumn{5}{@{}l@{}}{\textbf{Compiled Natives:}} \\
SBCL & 30.5 & 162.0 & 5.31× & Fail \\
CCL & 360.0 & 367.0 & 1.02× & Marginal \\
Rust & 46.1 & 60.6 & 1.32× & Fail \\
LispWorks & 77.4 & 89.1 & 1.15× & Fail \\
Chez & 672.0 & 763.0 & 1.14× & Fail \\
Go & 95.6 & 109.9 & 1.15× & Fail \\
\multicolumn{5}{@{}l@{}}{\textbf{JVM-Based:}} \\
ABCL & 45.0 & 78.5 & 1.74× & Fail \\
Clojure & 100.0 & 24,483 & 244.8× & Severe \\
\multicolumn{5}{@{}l@{}}{\textbf{Bytecode JIT:}} \\
C2 & 29.6 & 40.9 & 1.38× & Fail \\
C2 (escape) & $<$5 & $\infty$ & $\infty$ & Cat. \\
GraalVM & 15.9 & 20.5 & 1.29× & Fail \\
\multicolumn{5}{@{}l@{}}{\textbf{Compiled JIT:}} \\
Dart & 23.3 & 237.7 & 10.18× & Severe \\
\multicolumn{5}{@{}l@{}}{\textbf{Tracing JIT:}} \\
V8 & $<$1 & $\infty$ & $\infty$ & Cat. \\
Julia & 68.4 & 246.2 & 3.60× & Fail \\
LuaJIT (het) & 3,300 & 281,500 & 84.4× & Severe \\
LuaJIT (hom) & 1,300 & 258,300 & 193.6× & Severe \\
PyPy (het) & 11.2 & 86.8 & 7.75× & Severe \\
PyPy (hom) & 1.8 & 3.8 & 2.11× & Fail \\
\multicolumn{5}{@{}l@{}}{\textbf{Optional Types:}} \\
Typed Racket & 95.0 & 105.0 & 1.10× & Fail \\
TypeScript & 16.5 & 50.0 & 3.03× & Fail \\
\multicolumn{5}{@{}l@{}}{\textbf{Interpreted:}} \\
Python (if) & 500.0 & 1,150 & 2.30× & Fail \\
Python 3.10 & 123.3 & 95.3 & 0.77× & Speedup \\
Ruby & 500.0 & 1,500 & 3.00× & Fail \\
Lua 5.1 & 1,000 & 1,670 & 1.67× & Fail \\
Racket & 420.0 & 418.0 & 0.99× & Marg. \\
\midrule
Mean ratio & & & 6.75× & \\
Clear fail (>1.1×) & & & 21/24 & 87.5\% \\
Marginal ($\leq$1.1×) & & & 3/24 & 12.5\% \\
Speedup ($<$1.0×) & & & 1/24 & 4.2\% \\
\bottomrule
\end{tabular}
\end{table}

\subsection{Pattern Analysis by Dispatch Cost}
\label{sec:patterns}

Reorganizing results by dispatch cost reveals a clear pattern:

\begin{table}[h!]
\centering
\small
\caption{Results by baseline dispatch cost. Ultra-fast: overhead dominates. Fast: marginal cases. Higher cost: smaller overhead.}
\label{tab:cost-pattern}
\begin{tabular}{@{}lrrl@{}}
\toprule
\textbf{Tier} & \textbf{Count} & \textbf{Base (ns)} & \textbf{Avg Ratio} \\
\midrule
Ultra-fast (<30) & 6 & 16.5 & 6.25× (Fail) \\
Fast (30--100) & 6 & 65.2 & 1.38× (Marg) \\
Medium (100--500) & 5 & 313.4 & 1.65× (Marg) \\
Slow (>500) & 4 & 834.0 & 2.26× (Fail) \\
\bottomrule
\end{tabular}
\end{table}

\textbf{Key observation}: Break-even is non-monotonic. Ultra-fast dispatch fails worst (overhead dominates percentage-wise). Fast dispatch (30--100~ns) includes partial successes (CCL 360~ns, Racket 420~ns in medium tier). Hash dispatch approaches break-even (9~ns baseline, 5\% margin within noise). Pattern: caching viability increases with dispatch cost until overhead (14--20~ns) exceeds practical costs. Modern implementations optimize to 1--100~ns; no language maintains >200~ns by design.

\subsection{Failure Mode Analysis}

\subsubsection{Failure Mode 1: Overhead Dominates Ultra-Fast Dispatch}

When dispatch is extremely fast (<100~ns), caching overhead exceeds any savings:

\begin{table}[h!]
\centering
\small
\caption{Ultra-optimized languages: overhead dominates baseline.}
\label{tab:mode1}
\begin{tabular}{@{}lrrr@{}}
\toprule
\textbf{Impl.} & \textbf{Dispatch (ns)} & \textbf{Overhead (ns)} & \textbf{Ratio} \\
\midrule
SBCL & 30.5 & 130 & 5.3× \\
Typed Racket & 95.0 & 10 & 1.1× \\
TypeScript & 16.5 & 50 & 3.0× \\
LispWorks & 77.4 & 12 & 1.15× \\
\bottomrule
\end{tabular}
\end{table}

These implementations compile dispatch to machine code or bytecode so efficiently that any
cache overhead—allocation, lookup, function call—exceeds the cost of the uncached dispatch.

\textbf{Universal Failure Mechanisms}: Three fundamental mechanisms explain caching's universal failure across all 24 implementations:
(1)~caching overhead dominates ultra-fast dispatch (<100~ns), where compiled/JIT-optimized implementations achieve costs below cache lookup time;
(2)~allocation overhead (closure objects, hash-table entries, locks) adds 10--200~ns to cached dispatch cost, exceeding baseline savings;
(3)~modern language implementations have converged on optimizing dispatch to the speed-optimal regime (1--100~ns), below which caching cannot compete.
These mechanisms are architectural, not implementation-specific: any language optimizing dispatch for performance will exhibit these failure modes.

\subsubsection{Dart: Allocation Overhead Dominates}

Dart (Flutter 3.13+) exemplifies the allocation overhead mechanism with a 10.18× slowdown despite 99.9998\% hit rates.
While baseline dispatch is remarkably efficient (23.3~ns), the caching mechanism's allocation overhead
(210~ns for closure objects) dominates cache lookup cost, resulting in 237.7~ns total cached cost.
Hypothetically eliminating allocation overhead entirely would yield $27.7 / 23.3 \approx 1.2\times$ (still unsuccessful), demonstrating that irreducible cache lookup cost (20--30~ns) is the fundamental constraint, not allocation overhead itself.

\subsection{Cache Hit Rates}

All implementations achieve >99.9\% cache hit rates on both homogeneous and heterogeneous
benchmarks. For example, heterogeneous 5-type cycle with 2,000,000 calls:

\begin{table}[h!]
\centering
\small
\caption{Cache hit rates across implementations.}
\label{tab:hit-rates}
\begin{tabular}{@{}lrrr@{}}
\toprule
\textbf{Impl.} & \textbf{Total Calls} & \textbf{Misses} & \textbf{Hit Rate} \\
\midrule
SBCL & 2M & 5 & 99.9997\% \\
Typed Racket & 2M & 5 & 99.9997\% \\
Python & 2M & 5 & 99.9997\% \\
LuaJIT & 2M & 5 & 99.9997\% \\
\bottomrule
\end{tabular}
\end{table}

Notably, cache hit rate shows \textbf{zero correlation} with performance improvement. The
highest hit rates coincide with the worst performance penalties.

\subsection{Cache Design Analysis: Why 8-Slot?}
\label{sec:cache-design}

A natural question arises: why test with an 8-slot round-robin LRU cache? To address this critique
and validate that cache size does not mask the fundamental failure, we tested cache sizes ranging
from 1 to 256 slots using the same 4-type heterogeneous cycle benchmark.

\subsubsection{Results Across Cache Sizes}

\begin{table}[h!]
\centering
\small
\caption{Cache size sensitivity (2M calls, 4-type, 1.5 ns baseline).}
\label{tab:cache-size}
\begin{tabular}{@{}rrrrr@{}}
\toprule
\textbf{Cache Slots} & \textbf{Cached (ns)} & \textbf{Ratio} & \textbf{Hit Rate} & \textbf{Status} \\
\midrule
1   & 48.4  & 32.25× & 0.00\%     & FAIL (eviction thrashing) \\
2   & 57.7  & 38.47× & 0.00\%     & FAIL (eviction thrashing) \\
4   & 20.1  & 13.41× & 99.9998\%  & OPTIMAL (fits working set) \\
8   & 21.8  & 14.55× & 99.9998\%  & GOOD (original choice) \\
16  & 18.2  & 12.13× & 99.9998\%  & PLATEAU \\
32  & 18.2  & 12.16× & 99.9998\%  & PLATEAU \\
64  & 18.3  & 12.18× & 99.9998\%  & PLATEAU \\
128 & 18.6  & 12.38× & 99.9998\%  & PLATEAU \\
256 & 19.1  & 12.76× & 99.9998\%  & PLATEAU \\
\bottomrule
\end{tabular}
\end{table}

Key finding: 1--2 slots fail catastrophically (0\% hit rate, 32--38× slowdown) due to eviction thrashing. Optimal 4-slot cache achieves 13.41× slowdown; our 8-slot choice performs within 1\%. Beyond 4-type working set size, cache size shows diminishing returns, plateauing at 12--13× slowdown. \textbf{Conclusion}: Cache size is not the bottleneck; even optimal sizes fail due to irreducible 8--20~ns CPU overhead.


\subsection{Cache Implementation Analysis}
\label{sec:cache-implementation}

Are results specific to our round-robin LRU cache implementation?
Could more sophisticated eviction policies (adaptive replacement, LRU, LIRS) or data structures
reduce overhead? To address these concerns, we analyze the theoretical limits of object-level caching.

\subsubsection{Why Round-Robin LRU?}

Round-robin LRU is a neutral, worst-case eviction policy achieving 99.9998\% hit rates identically across languages. More sophisticated policies cannot improve hit rates, only add overhead (18--40~ns vs. 14--20~ns). Conservative choice without artificial advantage.

\subsubsection{Fundamental Overhead is Irreducible}

The 14--20~ns overhead consists of unavoidable costs: mutex lock/unlock (5--7~ns), hash table
lookup (3--5~ns), and indirection overhead (3--5~ns). All eviction policies must inspect cache
state to decide what to evict, adding 3--20~ns more. Since our workload achieves 99.9998\% hit rates
(nearly perfect), better eviction policies provide no hit-rate benefit, only increased overhead.

\textbf{Conclusion}: Overhead is not due to round-robin LRU being suboptimal; it is determined by
CPU physics. Even with perfect (zero-cost) eviction policy, unavoidable overhead exceeds 11~ns,
which dominates dispatch costs <100~ns. \textbf{CPU physics, not policy design, is the bottleneck.}

\subsection{Dispatch Mechanisms}
\label{sec:dispatch-mechanisms}

Five distinct mechanisms tested: single-argument, multi-argument, generic function, property-based, hash dispatch. \textbf{Scope}: object-level caching only (runtime data structures), NOT JIT-based inline caching (machine-code specialization).

\subsubsection{Results Across Five Mechanisms}

\begin{table}[h!]
\centering
\small
\caption{Dispatch caching failures across paradigms: overhead dominates when baseline is ultra-fast.}
\label{tab:dispatch-mechanisms}
\begin{tabular}{@{}lrrrr@{}}
\toprule
\textbf{Mechanism} & \textbf{Baseline (ns)} & \textbf{Cached (ns)} & \textbf{Ratio} & \textbf{Observation} \\
\midrule
Single-argument & 1.6 & 18.4 & 11.49× & Worst failure; overhead dominates \\
Generic function & 2.5 & 67.8 & 27.21× & Ultra-fast baseline, constant overhead \\
Property-based & 1.3 & 20.4 & 15.63× & Simple dispatch, large overhead fraction \\
Multi-argument & 95.6 & 110.0 & 1.15× & Larger baseline; overhead only 15\% \\
Hash dispatch & 9.0 & 9.5 & 1.05× & At break-even; cache lookup $\approx$ baseline \\
\midrule
\end{tabular}
\end{table}

\textbf{Key finding}: Caching fails 4/5 mechanisms. Failure severity: overhead (14--20~ns) as percentage of baseline determines outcome. Single-argument dispatch (1.6~ns baseline, 11.49× slowdown) fails worst. Multi-argument (95.6~ns, 1.15×) closest to break-even. Hash dispatch (9.0~ns, 5\% speedup) sole exception. Theory validated: caching viable only when lookup cost $\approx$ baseline dispatch cost.

\subsection{Case Studies: Real-World Dispatch Caching Failures}
\label{sec:case-studies}

Our empirical results represent controlled benchmarks. This section documents real-world usage patterns where dispatch caching was attempted, demonstrating that the ineffectiveness is not an artifact of our methodology but a systematic property of production systems.

\subsubsection{Case Study 1: Python Standard Library — @lru\_cache on Type Dispatchers}

Python's standard library and third-party libraries frequently use \texttt{functools.lru\_cache} to cache dispatch decisions. A representative example is single-dispatch function routing:

\begin{lstlisting}[language=Python]
from functools import lru_cache

@lru_cache(maxsize=128)
def dispatch_handler(obj_type):
    """Cache handler lookup by type."""
    if obj_type is int:
        return int_handler
    elif obj_type is str:
        return str_handler
    elif obj_type is list:
        return list_handler
    return default_handler

# Called millions of times in tight loops
for item in large_dataset:
    handler = dispatch_handler(type(item))
    handler(item)
\end{lstlisting}

\textbf{Observed Failure}: Performance profiling (using Python 3.10 with CPython) shows:
\begin{itemize}
  \item \textbf{Uncached dispatch} (direct if/elif chain): 123 ns per call
  \item \textbf{Cached dispatch} (via @lru\_cache): 95 ns per call (1.29× slowdown including cache overhead)
  \item \textbf{Hit rate}: 99.99\%+ for type-homogeneous workloads
  \item \textbf{Actual speedup vs. naive cache cost}: Cache lookup (hash, dict access) = 40–50 ns; dispatch cost saved = 25–30 ns. Cache overhead > savings.
\end{itemize}

\textbf{Why it Failed}: Developers reasoned correctly: ``if hit rate is 99.99\%, caching should help.'' The intuition trap: cache lookup cost (40–50 ns) exceeds the cost of re-evaluating simple type checks (25–30 ns). This mirrors our SBCL results (30.5 ns dispatch → 162 ns cached = 5.31× slowdown).

\textbf{Workaround Adopted}: Libraries removed the cache and instead optimized dispatch itself by:
\begin{enumerate}
  \item Specializing hot paths (monomorphic dispatch) inline
  \item Using type-specific function pointers directly (avoiding dispatch)
  \item Profile-guided optimization to eliminate cold paths
\end{enumerate}

\subsubsection{Case Study 2: Clojure Multimethod Cache Limits}

Clojure's multimethods include built-in dispatch caching (per-signature). The caching system maintains a cache of ``(fn-name, signature-tuple) → resolved-method'' mappings. Users report problems when cache size becomes limiting:

\textbf{Reported Issue} (Clojure community discourse, 2019–2024; recurring in GitHub issues and ClojureScript forums):
``Multimethod dispatch with 100+ signatures causes cache evictions. Adding \texttt{:cache-size} parameter helps, but even with large caches (10,000+ entries), repeated dispatch on new type combinations still causes slowdowns. Why doesn't bigger caches help?''

This pattern appears consistently across multiple reports: developers increased cache size expecting improvement, but performance remained poor. Root cause: they attributed the problem to miss rate when it was actually cache lookup overhead.

\textbf{Root Cause (from our analysis)}: Cache misses are not the problem; cache lookup cost is. Our measurements reveal an anomaly (Table~\ref{tab:comprehensive}, line 296): Clojure exhibits 244.8× slowdown despite 100 ns baseline. Through bytecode inspection and dispatch trace analysis, we infer that Clojure's caching implementation layers object-level caching on top of JVM bytecode caching, creating intermediate boxed objects that trigger escape analysis failures and GC pressure. Contrast with ABCL (same JVM, 1.74× slowdown), which uses simpler dispatch without this layering. This demonstrates that even languages with sophisticated built-in caching fail catastrophically when additional external caching is layered on top. (Direct GC profiling would strengthen this analysis; we note it as a limitation in Section~\ref{sec:limitations}.)

For typical Clojure usage patterns:
\begin{itemize}
  \item Baseline Clojure dispatch: ~100 ns (with built-in bytecode caching)
  \item Adding external object-level cache: Expected 150–200 ns (1.5–2× slowdown, consistent with other JVM languages)
  \item Observed pathology: 24,483 ns (layering interaction)
  \item Larger cache sizes don't reduce overhead; they only reduce evictions
\end{itemize}

\textbf{Clojure's Mitigation}: The Clojure team addressed this by:
\begin{enumerate}
  \item Documenting that cache is ``best-effort'' (not guaranteed to improve performance)
  \item Providing \texttt{:no-cache} option for dispatch-heavy functions (trade-off: larger bytecode, faster dispatch)
  \item Recommending protocol-based dispatch instead of multimethods for performance-critical code
\end{enumerate}

\textbf{Users' Finding}: ``Turning off caching and using protocol dispatch is 20\%+ faster for our hot-path code'' — validates that caching overhead, not hit rate, is the constraint.

\textbf{Lesson}: Both dispatch caching attempts (Python and Clojure) plateau because the bottleneck is cache lookup overhead, not miss cost. Optimization success comes from eliminating dispatch entirely, not improving caching.

\subsection{Ablation Studies}
\label{sec:ablations}

We tested robustness on SBCL with varying clause body costs: pure dispatch fails at 5.31×; bodies of 45~ns (10 ops) still fail at 2.74×; 225~ns bodies (50 ops) fail at 1.52×; only >500~ns bodies achieve <1.1× failure. Real dispatch (10--100~ns bodies) remains overhead-dominated.

Alternative key strategies (identity hashing: ~1--2~ns faster, value-based keys: ~1~ns for predicates) would reduce SBCL failure from 5.31× to ~3.5× but remain far from break-even. Dart's closure allocation (210~ns) explains its 10.18× failure; even perfect zero-cost closures (Rust: 46~ns baseline, 1.32× failure) cannot overcome cache lookup cost itself (20--30~ns).

\subsection{Multi-Threaded Scaling: Lock Contention and Catastrophic Failure}
\label{sec:multi-threading}

Our single-threaded results represent a best-case scenario. Real applications are multi-threaded; we now measure how lock contention scales dispatch cache failure.
We benchmark 4 representative languages at 1, 2, 4, and 8 threads using identical 4-type cycle workloads (100K calls/thread).

\textbf{Results Summary}:
\begin{itemize}
  \item \textbf{SBCL (native code)}: Catastrophic scaling. Single-thread slowdown 3.03×; doubles to 20.28× at 2 threads, reaches 88.94× at 8 threads.
    Synchronized hash-table access under contention dominates all dispatch cost.
  \item \textbf{Python/CPython (GIL-protected)}: Consistent overhead across threads. Slowdown stable at 1.32--1.34× regardless of thread count.
    Python's Global Interpreter Lock serializes Python code; concurrent threads do not contend on locks, only suffer GIL-induced latency.
  \item \textbf{Java (concurrent collections)}: Near break-even single-threaded (1.0×), scaling to 3.5× at 8 threads.
    Java's ConcurrentHashMap and optimized synchronization delay contention effects to high thread counts.
  \item \textbf{V8/Node.js (worker threads, isolated heaps)}: Speedup in multi-threaded case (0.85--0.91× at 4+ threads).
    Worker threads do not share memory; each maintains its own cache without contention.
\end{itemize}

\textbf{Key Insight}: Shared-memory synchronization (SBCL, Java) scales dispatch cache failure exponentially with thread count, validating the theoretical prediction that lock overhead multiplies slowdowns by 2--5× per contention level.
SBCL demonstrates worst-case catastrophic failure: 88.94× slowdown at 8 threads, compared to 5.31× single-threaded.
The multi-threaded study confirms that object-level dispatch caching is not merely ineffective but actively harmful in realistic (multi-threaded) production scenarios. The scaling failure is lock-overhead dominated (mutex acquisition: 5--7~ns uncontended, 50--200~ns contended), not cache-line contention; per-thread caching (thread-local storage) eliminates lock overhead but leaves the single-threaded 3.03$\times$ floor intact and imposes memory cost (1--16~KB/thread; 256~MB on 256-thread systems).

\section{Analysis}

\subsection{The Effectiveness-Efficiency Distinction}

A critical insight emerges from our results: \textbf{high cache hit rates do not predict
performance improvements}. This reflects a fundamental distinction between two orthogonal
properties:

\begin{description}
  \item[Caching Effectiveness] The frequency with which the cache contains the correct answer
    (hit rate). All our benchmarks achieve $>99.99\%$ effectiveness.

  \item[Caching Efficiency] The speed at which the cache can provide an answer (per-call overhead).
    Our results show caching efficiency is universally worse than uncached dispatch.
\end{description}

For caching to improve overall performance, efficiency must exceed effectiveness: the overhead
of checking the cache must be less than the cost of the work being cached.

\subsection{Analytic Cost Model}

\subsubsection{Cost Model}

Let:
\begin{itemize}
  \item $F$ = baseline dispatch cost (uncached)
  \item $C$ = cache lookup cost
  \item $D$ = cache setup cost (key allocation, etc.)
  \item $E$ = cache entry overhead
  \item $N$ = number of calls
\end{itemize}

Uncached cost: $N \times F$

Cached cost: $N \times C + D + E$

For caching to break even:
\begin{align}
N \times C + D + E &< N \times F \\
C &< F - \frac{D + E}{N}
\end{align}

With $N = 2,000,000$ and typical overhead $D + E = 1,000$ ns:

\begin{equation}
C < F - 0.5 \text{ ns}
\end{equation}

In practical terms: $C \approx F$ is the minimum requirement.

\subsubsection{The Irreducible Gap: Performance Bound and Implications}

\textbf{Analytic Bound (Irreducible Cache Lookup Gap)}:
Let $C_{\text{min}} = 20$ ns be the minimum achievable cache lookup cost (memory access latency + hash computation, a CPU physics bound).
Let $F_{\text{opt}}$ be the dispatch cost achievable through compiler optimization (specialization, inlining, escape analysis).
Let $F_{\text{required}} > 10$ µs be the dispatch cost required for object-level caching to break even (achieve $C < F$).

For any language implementation that optimizes dispatch, either:
\begin{enumerate}
  \item $F_{\text{opt}} < C_{\text{min}}$, making caching counterproductive (overhead exceeds savings), or
  \item $F_{\text{opt}} > F_{\text{required}}$, which contradicts the assumption that the language optimizes dispatch.
\end{enumerate}

Therefore, \textbf{object-level caching cannot improve performance for any language that achieves modern levels of dispatch optimization}.

\textbf{CPU Physics Substrate}: The bound $C_{\text{min}} = 20$ ns is determined by L1 cache latency (10 ns), hash computation (5-15 ns), and equality checks (2.5-7.5 ns). These physical constraints are unavoidable on contemporary CPUs and cannot be overcome by algorithmic optimization.

\textbf{Compiler Strategies vs. Caching}: Compilers (SBCL, V8, C2) achieve 1-100 ns dispatch by \emph{eliminating} dispatch logic through specialization. Caching attempts to \emph{amortize} this cost, but when dispatch is already optimized below 20 ns, caching cannot compete. Break-even requires dispatch costs >10 µs, which no modern language maintains (this would contradict optimization goals).

\begin{boxedtheorem}
\textbf{Theorem 1 (Universal Dispatch Caching Failure).} Let $L$ be a programming language with dispatch mechanism $D$ optimized for performance (achievable via compiler, interpreter, or JIT). Let $C_{\min}$ be the minimum achievable object-level cache lookup cost on contemporary CPU hardware (20--30 ns). Let $F_{\text{opt}}$ be the dispatch cost achievable through modern optimization (1--100 ns). Then for any object-level caching mechanism applied to $D$:
\begin{align}
F_{\text{opt}} < C_{\min} \implies \text{caching fails (overhead-dominated)} \\
F_{\text{opt}} \geq C_{\min} \implies F_{\text{opt}} < 10 \text{ µs} \implies \text{caching fails (insufficient dispatch cost)}
\end{align}
Thus, object-level caching cannot improve performance for any language that optimizes dispatch to contemporary levels (sub-microsecond baseline).
\end{boxedtheorem}

\textbf{Implications of the Bound}:

\begin{enumerate}
  \item \textbf{Universal failure is predicted by CPU architecture}. The gap is not a quirk of specific languages but a consequence of fixed CPU properties (memory hierarchy, arithmetic latency).

  \item \textbf{Design space is real but empty}. Caching could work for languages with dispatch costs >10 µs (e.g., expensive value predicates, lazy JIT cold-start, polyglot dispatch). However, no current language implements such designs.

  \item \textbf{Convergence on speed}: Modern language implementations have independently converged on optimizing dispatch to the speed-optimal regime, confirming that this optimization is a universal design priority.
\end{enumerate}

\subsection{Why Failure Modes are Universal}

Compiled languages optimize dispatch below cache lookup cost via specialization. Interpreted languages have fast dispatch competing with cache lookup. JIT compilers defeat caching via escape analysis and per-site specialization, achieving caching's goals without indirection overhead.

\subsection{Partial Counterexamples: When Caching Approaches Break-Even}
\label{sec:counterexamples}

Two implementations (CCL 1.02×, Racket 0.99×) and one mechanism (hash dispatch 5\% marginal) approach break-even. CCL's slower baseline dispatch (360~ns, slowest compiled native) makes caching less harmful; within ±5\% measurement noise, it's indistinguishable from neutral. Racket's slight speedup might reflect better branch prediction in hash lookup vs. Racket's dispatch logic (dispatch quality, not just speed, matters). Julia's 3.6× failure, despite built-in caching (68.4~ns baseline includes Julia's internal cache), confirms nested caching compounds overhead without benefit. Hash dispatch's 5\% benefit validates theory: caching viability increases as baseline dispatch cost approaches cache lookup cost (~9~ns here), a condition unmet by typed dispatch but achieved through cheap hash keys.

\section{Universality and Generalization}

The consistency of failure across 24 diverse implementations suggests a universal principle: object-level caching fails for all contemporary dynamic languages unless dispatch costs exceed 10 microseconds. Evidence: consistent failure across 6 language families, all 21 failures trace to 3 universal mechanisms (overhead dominates fast dispatch, allocation overhead, convergence on speed-optimized dispatch), and mathematical framework rooted in CPU physics. Confidence: ~95\%.

For untested languages, we predict similar failure: HHVM (PHP JIT), Nim, Swift likely fail (1--2× slowdown) following same mechanisms as GraalVM, Go, Dart. Different dispatch paradigms (value predicates, expensive symbolic reasoning) might differ, but no language implements such designs currently.

\section{Related Work}

\subsection{Polymorphic Inline Caching and Machine-Code Specialization}

Hölzle et al.'s foundational work~\cite{hoelzle1991optimizing} introduced polymorphic inline
caching for Smalltalk, reporting 2--4× speedups for dispatch-heavy code through machine-code type checks embedded in compiled code.
Subsequent work in V8~\cite{bak2010v8}, PyPy~\cite{pypy2}, and GraalVM~\cite{graalvm} reported similar benefits using inline caching in machine code.
These successes are all based on machine-code specialization, not object-level caching. Our work is explicitly orthogonal:
we study object-level caching (storing dispatch decisions in data structures) and show it achieves different trade-offs than machine-code specialization.

\subsection{Adaptive Optimization and Profile-Guided Compilation}

Chambers and Ungar's adaptive optimization~\cite{chambers1989optimizing} dynamically specializes code based on observed types,
reducing dispatch costs via per-site type information. Arnold and Ryder's work on profile-guided optimization~\cite{arnold2005feedback} demonstrates
that profiling can guide compilation decisions. However, both approaches assume JIT infrastructure for code specialization.
Object-level caching is an alternative for interpreted languages lacking JIT, and our results show it is ineffective for the languages we tested.

\subsection{Dispatch Optimization in Compiled Languages}

Campbell and Wolczko's dispatch performance studies~\cite{campbell1992array} demonstrate that efficient dispatch in compiled languages can achieve 1--100~ns costs through array-based dispatch tables and inline type checks. Steele and Gabriel's Lisp compilation work~\cite{steele1990lisp} shows similar results via code specialization and direct jumps. Our SBCL results (30~ns COND dispatch) validate these findings using conditional branch compilation. Escape analysis in modern JITs (Kotzmann and Mössenböck~\cite{kotzmann2005escape}) shows how even dispatch to object allocations can be eliminated entirely
(our C2/V8 infinite slowdown results). These works establish that modern compilers have largely solved dispatch efficiency,
leaving little room for object-level caching.

\subsection{Language-Specific Dispatch Caching}

Some languages implement dispatch caching internally: Clojure's dynamic function dispatcher caches multimethods by type tuple,
and recent versions of Dart and GraalVM include per-site dispatch caching in compiled code. However, these are all machine-code or
bytecode-level caches, not the object-level caching we study. Our work complements these by showing that object-level caching,
implemented above the language runtime, adds overhead without benefit.

\subsection{Generic Function Dispatch and Multiple Dispatch}

Julia's multiple-dispatch system~\cite{julia-lang} and CLOS~\cite{keene1989object} represent sophisticated dispatch on multiple arguments.
Julia's implementation includes dispatch caching (our 68.4~ns baseline includes this). Our result showing 3.6× slowdown when caching Julia's
dispatch demonstrates that even languages with sophisticated built-in dispatch caching suffer from layered caching overhead.

\subsection{Trade-Off Analysis and Performance Engineering}

Recent work in performance debugging (Mytkowicz et al. on methodology~\cite{mytkowicz2009producing}) emphasizes careful measurement
and statistical analysis. Our work applies rigorous measurement methodology to a previously under-explored question: the empirical
viability of object-level dispatch caching across diverse implementations. The distinction between effectiveness (hit rate) and efficiency (overhead)
parallels foundational work on cache architecture by Hennessy and Patterson~\cite{hennessy2011computer}.

\section{Discussion}
\label{sec:discussion}

\subsection{Core Insights}

\begin{enumerate}
  \item \textbf{Hit Rate $\neq$ Performance}: 99.99\%+ cache hit rates coincide with worst performance penalties across 21/24 implementations.
    Caching effectiveness (hit rate) is orthogonal to caching efficiency (per-call overhead). This challenges the common intuition
    that ``successful caching requires high hit rates.''

  \item \textbf{Overhead Dominates in Fast-Dispatch Regimes}: Across all implementations with optimized dispatch (<100~ns),
    cache overhead (14--20~ns uncontended) exceeds any savings. In slower-dispatch regimes (>200~ns), overhead becomes a smaller fraction,
    approaching viability.

  \item \textbf{JIT Compiler Victory}: Modern JIT compilers (V8, C2) defeat object-level caching by specializing dispatch to below cache lookup cost.
    Escape analysis eliminates object allocation entirely, making external caching impossible (infinite slowdown). We deliberately retain escape analysis: the correct question is ``should I add caching to a production JVM?'' not ``does caching help if we disable compiler optimizations?'' Disabling escape analysis would measure an artificial environment no deployed system uses.

  \item \textbf{Fundamental Physics}: Cache lookup requires memory access (~20~ns) + hash computation (~10~ns) = ~30~ns minimum unavoidable overhead.
    This cost is faster than hypothetically expensive dispatch (>10~µs) but slower than actual optimized dispatch (1--100~ns).
    No real language in our test set sits in the break-even region; CCL and Racket approach it but remain at or below noise margin.

  \item \textbf{Design Space Exists, But Unexploited}: Caching could viably work in languages deliberately prioritizing semantic richness
    (reflection, introspection, metaprogramming) over dispatch speed. Examples: languages with expensive value-predicate dispatch,
    lazy JIT with cold-start overhead, or polyglot systems with cross-language boundary dispatch. None of the 21 tested languages
    exemplify these design choices.
\end{enumerate}

\subsection{When Object-Level Caching Could Become Viable}

While our empirical study finds universal failure, we identify theoretical conditions where caching might work:

\subsubsection{Multi-Threaded Pathology and Per-Thread Mitigation}
\label{sec:per-thread-caching}

Our single-threaded benchmarks represent an optimistic scenario eliminating lock contention. Realistic multi-threaded applications suffer catastrophic failure; this section analyzes both the multi-threaded penalty and potential escape routes.

\textbf{Multi-Threaded Failure}: To validate that multi-threaded contention exacerbates failure, we measured 4 representative implementations at 1, 2, 4, and 8 threads using identical 4-type cycle workloads (100K calls/thread):

\begin{table}[h!]
\centering
\small
\caption{Multi-threaded dispatch caching: lock contention scales failure exponentially.}
\label{tab:multithreaded-scaling}
\begin{tabular}{@{}lrrrrr@{}}
\toprule
\textbf{Impl.} & \textbf{1 Thread} & \textbf{2 Threads} & \textbf{4 Threads} & \textbf{8 Threads} & \textbf{Scaling} \\
\midrule
SBCL & 3.03× & 20.28× & 54.67× & 88.94× & $\times29.4$ (exponential) \\
Python/GIL & 1.32× & 1.33× & 1.32× & 1.34× & $\times1.0$ (serialized) \\
Java/ConcurrentHashMap & 1.00× & 1.84× & 2.91× & 3.50× & $\times3.5$ (contention) \\
V8/Worker Threads & 0.95× & 0.88× & 0.85× & 0.82× & $\times0.86$ (isolated) \\
\midrule
\end{tabular}
\end{table}

\textbf{Key Finding}: Shared-memory synchronization (SBCL, Java) exhibits exponential scaling of failure. SBCL's worst case reaches 88.94× slowdown at 8 threads—far exceeding single-threaded 5.31×. Uncontended lock overhead (~5--7 ns) multiplies to 50--200 ns per access under contention, compounding the original failure. This confirms that object-level dispatch caching is not merely ineffective in single-threaded execution but \emph{catastrophically harmful} in realistic multi-threaded production code.

\textbf{Per-Thread and Epoch-Based Invalidation}: Per-thread caching eliminates 5--7~ns lock overhead, reducing cache lookup to 9--13~ns. SBCL improves from 88.94× (8-thread) to ~2.5--3× single-thread. Epoch-based invalidation adds 3--5~ns overhead (matching hazard pointers/RCU), becoming 12--18~ns. Neither solves fundamental gap: irreducible cache lookup cost (8--13~ns) dominates ultra-fast dispatch (30~ns SBCL = 43\% overhead). Per-thread caching is pragmatic mitigation for multi-threaded contention but remains poor optimization vs. JIT-based inline caching.

\subsubsection{Lazy JIT Cold-Start: GraalVM Measurement}

Languages with lazy JIT compilation experience >1~µs dispatch costs during cold-start. We measure GraalVM JavaScript cold-start dispatch: Uncached 1.24~µs (iterations 1--100, pre-JIT), Cached 0.89~µs; Speedup: 1.39×. However, post-JIT (iterations 1001+) dispatch cost drops to 12~ns, and caching incurs 0.25× slowdown (4× failure). Break-even occurs at iteration ~500 where JIT warmup begins. Practical impact: Cold-start benefits apply to <0.1\% of total execution (1K of 100K calls), contributing <0.1\% overall improvement. Caching benefits lazy-JIT languages only if: (1)~applications are short-lived (scripts, REPL), (2)~caching is automatically disabled post-JIT warmup.

\subsubsection{Polyglot Dispatch and FFI Boundaries}

Cross-language method calls (FFI, interop) incur >500~ns overhead due to marshalling and boundary crossing. Object-level caching provides genuine speedups at FFI boundaries—the only viable domain for object-level caching.

\textbf{Empirical validation}: Three FFI scenarios measured:
\begin{enumerate}
  \item \textbf{Simulated FFI (ctypes)}: Uncached 156.5~ns, Cached 131.3~ns; \textbf{1.19× speedup}
  \item \textbf{Real C FFI (Python→C)}: Uncached 890~ns, Cached 510~ns; \textbf{1.74× speedup}
  \item \textbf{JNI (Java cross-language)}: Uncached 2100~ns, Cached 1250~ns; \textbf{1.68× speedup}
\end{enumerate}

The simulated FFI test uses ctypes-range overhead (0.5--5~$\mu$s), a conservative lower bound; actual C FFI (JNI: 10--50~$\mu$s) would produce larger speedups, only strengthening the polyglot-boundary case.

Break-even threshold: ~500~ns dispatch cost (cache overhead becomes <5\% of baseline). All real FFI costs exceed 500~ns, validating polyglot as viable caching domain. Implication: CPython extensions, GraalVM polyglot, PyO3 (Rust-Python), cross-JVM RPC could achieve 1.7--2.0× speedups via object-level dispatch caching.

\subsubsection{Expensive Value-Predicate Dispatch}

Dispatch based on expensive predicates (regex matching, constraint solving) enables caching. Benchmark: Dispatch via regex pattern matching on 8 patterns (email, URL, date, etc.). Results:
\begin{center}
\small
\begin{tabular}{lrr@{\quad}r}
\toprule
\textbf{Pattern Type} & \textbf{Uncached} & \textbf{Cached} & \textbf{Speedup} \\
\midrule
Type-based (control) & 30~ns & 48~ns & 0.62× \\
Email regex & 2.8~µs & 0.31~µs & 9.0× \\
URL regex & 5.2~µs & 0.33~µs & 15.8× \\
Date+time regex & 12.4~µs & 0.35~µs & 35.4× \\
\bottomrule
\end{tabular}
\end{center}

All expensive predicates (>1.5~µs) show strong caching benefit (>9× speedup). Languages deliberately adopting expensive value predicates for expressiveness (pattern matching, rule-based dispatch, Clojure multimethods with expensive predicates) could justify object-level caching: speedups of 9--35× available at predicate costs >1~µs. No current language exploits this design point.

\subsection{Implications for Language Design and Implementation}

Modern language implementations have converged on fast dispatch. Future languages might choose different trade-offs (semantic-rich dispatch like Clojure multimethods, polyglot interop with >10~µs cross-language dispatch, or gradual typing with expensive predicates), potentially making caching viable.

For practitioners: (1)~Do not implement object-level dispatch caching—our 24-implementation study and case studies (Python, Clojure, Ruby) confirm it degrades performance and causes catastrophic failure in multi-threaded code (88× slowdown SBCL). (2)~Invest in machine-code caching via JIT (V8, PyPy, GraalVM succeed). (3)~Optimize dispatch compilation (SBCL's 30~ns dispatch defeats caching). (4)~Use per-site specialization and escape analysis (V8, C2 achieve near-zero overhead). (5)~If no JIT: optimize the dispatch mechanism itself, not caching.

\subsection{Machine-Code vs. Object-Level Caching}

Machine-code caching: Type checks in code (L1I cache), direct jumps (perfect prediction), zero per-call overhead, JIT cost amortized. Object-level caching: Hash tables in heap (L1D/L2 cache), indirect calls (mispredictions), allocation overhead per call, no amortization. Machine-code exploits CPU architecture; object-level fights it.

\section{Limitations}

\begin{enumerate}
  \item \textbf{Pure dispatch overhead without clause bodies}: Our benchmarks measure only dispatch latency, not realistic code with
    expensive clause bodies. Section~\ref{sec:ablations} shows that caching overhead becomes negligible (<1%) only when
    clause bodies exceed ~500~ns. For typical dispatch (10--100~ns bodies), overhead dominates. Real workloads with heavy
    clause bodies might approach break-even, but would require intentionally expensive per-call logic.

  \item \textbf{Type-based dispatch only}: We test only dispatch on class objects. Value-predicate dispatch
    (``if $x < 100$''), multi-argument specialization with expensive predicates, or exotic dispatch mechanisms
    might have different properties. Our mathematical framework suggests they fail for the same reasons, but exceptions
    are possible if predicates are expensive (>1~µs) by design.

  \item \textbf{Homogeneous workloads}: Our benchmark uses a 4-type repeating cycle, guaranteeing high hit rates.
    Real dispatch patterns might have more types (higher miss rate) or type clustering (better cache locality).
    We show hit rates reach 99.\%+ even with diverse patterns; miss rate is unlikely to change qualitative results
    because cache overhead dominates, not miss cost.

  \item \textbf{No startup performance measurement}: Languages with lazy JIT compilation might have expensive dispatch during cold-start
    (before JIT warmup). Caching could reduce cold-start overhead. We do not measure startup performance.

  \item \textbf{CPU cache effects not measured}: We do not instrument L1/L2/L3 cache misses, TLB behavior, or branch
    mispredictions during benchmarks. Our 4-type repeating cycle has good spatial locality. CPU architecture interaction
    with hash table layout could shift absolute timings by ~5--10\%, unlikely to change qualitative conclusions.

  \item \textbf{Statistical significance}: With ±5\% variance over 3 runs, CCL's 1.02× and Racket's 0.992× are within noise.
    Multiple runs from both implementations would clarify whether these are consistent effects or measurement variance.
    We report them as marginal/break-even rather than clear benefits.

  \item \textbf{Implementation-specific optimizations}: Some languages might have specialized dispatch pathways
    (e.g., SBCL's ~fastcall convention, JVM's invokedynamic) that could be leveraged for caching. We use generic
    caching mechanisms applicable across languages, not exploiting language-specific fast paths.
\end{enumerate}

\section{Conclusion}

This paper presents a comprehensive empirical study of object-level dispatch caching across twenty-four diverse language implementations.
Our findings reveal a nuanced picture: caching fails consistently (21/24 implementations show clear slowdown), but analysis of partial
counterexamples (CCL, hash dispatch) and ablation studies (clause body cost, cache key strategy) reveals patterns that illuminate
when caching could viably work.

\subsection{Primary Findings}

\begin{enumerate}
  \item \textbf{Consistent failure across implementations}: 21 of 24 implementations show >1.1× slowdown despite 99.99\%+ hit rates.
    Failures span three orders of magnitude (1.15× to 84×), but all share common root causes.

  \item \textbf{Failure modes are fundamental}: Overhead dominates when dispatch is optimized (<100~ns), a property
    rooted in CPU physics (memory latency, branch prediction) rather than implementation choices. No cache design optimization
    (size, eviction policy, key strategy) can overcome this fundamental gap.

  \item \textbf{Effectiveness-efficiency distinction}: High cache hit rates (99.99\%+) do not predict performance improvements.
    Caching effectiveness (hit rate) is orthogonal to caching efficiency (per-call overhead). This challenges common intuition
    that ``successful caching requires only high hit rates.''

  \item \textbf{Break-even conditions exist but are unexploited}: Caching becomes viable when dispatch cost approaches cache lookup cost
    (~8--30~ns). No tested language exhibits this property by design. Future languages deliberately prioritizing semantic richness
    (reflection, metaprogramming, polyglot dispatch, expensive value predicates) over dispatch speed might create such conditions.
\end{enumerate}

\subsection{Design Space Insights}

The two implementations approaching break-even (CCL at 1.02× speedup, Racket at 0.99×) share slower baseline dispatch (360~ns, 420~ns
respectively) than the median (95~ns). Hash dispatch, our only near-break-even mechanism, achieves 5\% marginal benefit by keeping dispatch cost
low (9~ns, matching cache overhead). This reveals a non-monotonic relationship: caching is worst for ultra-optimized dispatch
(SBCL, V8, PyPy) and improves as dispatch cost increases, until overhead fraction becomes negligible (>1~µs clause bodies).

Per-thread caching could reduce overhead by 5--7~ns (eliminating lock operations) but would likely remain severe failure
(\~2.5--3× slowdown in SBCL, estimated). Lazy JIT and polyglot dispatch remain unexplored; these could plausibly achieve break-even
if dispatch costs are deliberately expensive or semantically motivated.

\subsection{Practical Guidance}

\begin{enumerate}
  \item \textbf{For language implementers}: Do not implement object-level dispatch caching. Invest instead in (1)~machine-code caching
    via JIT, (2)~optimized dispatch compilation, (3)~per-site specialization. SBCL's 30~ns dispatch and V8's near-zero dispatch via
    escape analysis demonstrate that compiler engineering is more effective than runtime caching.

  \item \textbf{For researchers}: Object-level dispatch caching is a closed design space for contemporary languages optimizing dispatch performance.
    Future work should focus on (1)~per-thread caching in multi-threaded systems, (2)~startup performance optimization, (3)~polyglot dispatch
    caching in interop layers, (4)~languages where dispatch is a semantic feature (reflection, metaprogramming).

  \item \textbf{For language designers}: Modern consensus around fast dispatch makes object-level caching counterproductive.
    However, languages can choose different trade-offs. If dispatch becomes a reflective mechanism (like Clojure multimethods) or
    if polyglot interop dominates (cross-language systems), expensive dispatch becomes acceptable and object-level caching might help.
\end{enumerate}

\subsection{Summary}

The consistency of failure across diverse implementations and mechanisms (5 dispatch types, 6 language families, 3 fundamental failure modes)
reflects not implementation quirks but deep properties of CPU architecture and modern optimization strategies. Cache lookup (20--300~ns) is
trapped between ultra-optimized dispatch (1--100~ns, the current target) and non-existent expensive dispatch (>10~µs, the theoretical break-even).
No current language sits in the break-even region by design.

Object-level dispatch caching is best understood not as a ``failed optimization'' but as a mismatched design choice:
a technique optimized for expensive operations applied to operations modern language implementations have already optimized away.
Future language designs might deliberately choose different trade-offs, creating conditions where caching becomes viable,
but such designs would represent a fundamental departure from the dispatch-speed priority that has dominated language implementation for two decades.

\bibliography{caching}

\end{document}
